\documentclass{beamer}
\usepackage[T1]{fontenc}
\usepackage{newpxtext,newpxmath}
\usepackage[varqu,varl,scale=.95]{inconsolata}
\usepackage{csquotes}
\usepackage[ghc,language=haskell,xlistings={style={plain number}}]{magic-haskell}
% a calm palette (after Catppuccin Latte) through the colors of xlistings
\definecolor{plum}{HTML}{9B1B6E}         \definecolor{latte-blue}{HTML}{1E66F5}
\definecolor{latte-green}{HTML}{40A02B}  \definecolor{latte-text}{HTML}{4C4F69}
\definecolor{latte-subtle}{HTML}{9CA0B0}
\lstcolorlet{keywordA}{plum}             % keywords
\lstcolorlet{keywordB}{latte-blue}       % types and constructors
\lstcolorlet{keywordC}{latte-text}       % functions, marked by their links
\lstcolorlet{numbers}{latte-blue}
\lstcolorlet{literals}{latte-green}
\lstcolorlet{comments}{latte-subtle}
\usefonttheme{serif}
\title{\enquote{Magic} Haskell}
\begin{document}
\begin{frame}[fragile]{A GHCi session, step by step}
Each step appears on its own overlay:
\begin{ghcisession}
let xs = [1 .. 10]
sum xs
filter even xs
:t foldr
foldr (\x acc -> x + acc) 0 xs
\end{ghcisession}
\end{frame}
\begin{frame}[fragile]{Types from GHC}
\begin{minted}{haskell}
fib :: Integer -> Integer
fib n = if n < 2 then n else fib (n - 1) + fib (n - 2)
\end{minted}
The type of \bhaskell{zipWith} is \mhsType{zipWith} and \bhaskell{take 5 [2, 4 ..]} is \mhsEval{take 5 [2, 4 ..]}.
\end{frame}
\end{document}
